% !TeX root = ../教材题库_全部条目_转写工作稿.tex
% Source: XJ-B1-C06-TH-P242-Y0889
% 第6章 / 6.4 用样本估计总体 / 教材 p.242、243 / PDF 第249、250页
% OCR未校对：公式与文字均可能有误；未完成数学TeX转写；禁止直接发给学生
% 原题及相邻原解截图（核对用）：../题目原图/XJ-B1-C06-TH-P242-Y0889-01.webp, ../题目原图/XJ-B1-C06-TH-P242-Y0889-02.webp
% 本文件只在转写完成后作为题干使用；OCR 条目尚不是可靠题干。
\begin{bankitem}{p.242 思考（含紧邻说明）}{第6章 / 6.4 用样本估计总体 / 教材 p.242、243 / PDF 第249、250页}{XJ-B1-C06-TH-P242-Y0889}{OCR未校对：公式与文字均可能有误；未完成数学TeX转写；禁止直接发给学生}
\begin{ocrdraft}
[教材p.242 / PDF第249页]
\par
？思考育才中学从甲、乙两名射击运动员中选拔一人参加市中学生运动会，
\par
甲、乙两人参加测试的成绩（单位：环）如下：
\par
甲：7，8，8，9，7，8，8，9，7，9；乙：6，8，7，
\par
7，8，9，10，
\par
7，9，9．
\par
你认为教练员该如何选出合适选手？
\par
首先考虑两人射击测试的平均成绩，经计算得《甲一\&0，》乙一\&0，可见两人
\par
的平均成绩相同．因此从两人的平均成绩角度无法选出合适选手．
\par
我们可以将甲、乙的射击成绩表示在图6．4一2中．
\par
甲射击成绩／环
\par
乙射击成绩/环
\par
10
\par
10
\par
9
\par
9
\par
8
\par
8
\par
7
\par
7
\par
6
\par
6
\par
5
\par
5
\par
2
\par
3
\par
4
\par
5
\par
67
\par
8
\par
910射击次数
\par
12345678910射击次数
\par
（2）
\par
图6．4一2
\par
比较上面两幅图可以发现，甲的射击成绩大多集中在平均成绩8环的附近，而
\par
乙的射击成绩与平均成绩比较，波动较大．
\par
如何刻画一组数据的波动大小或离散程度呢？
\par
[教材p.243 / PDF第250页]
\par
一种简单办法是计算极差．计算可得甲、乙两人10次射击成绩的极差分别为2
\par
和4，可以发现乙的成绩波动大些．但由于极差只利用了一组数据两端的信息，对
\par
其他数据的取值情况没有涉及，不能反映中间数据的离散状况，因而不能全面地刻
\par
画数据的离散程度．
\par
统计学中，常采用方差来刻画一组数据波动的大小．若设，3，2，
\par
\ensuremath{\cdot}，”是总
\par
体的全部个体，“是总体均值，则称
\par
（夢一萨）2十（3'2一）2十'\ensuremath{\cdot}十(YN\ocrglyph{00AF}g)2
\par
为总体方差．
\par
总体方差。2刻画了总体中的个体向总体均值“的集中或离
\par
散的程度：方差越小，表明个体与均值的距离越近，个体
\par
方差越大又表示什么？
\par
向“集中得越好．总体方差。2也刻画了总体中个体的稳定或
\par
波动的程度：方差越小，表明个体越整齐，波动越小．
\par
类似地，若从总体中随机抽样，获得”个观测数据
\par
，用》表示
\par
这”个数据的均值，则称
\par
1
\par
为这”个数据的样本方差．
\par
样本方差刻画了样本数据相对于样本均值集中或离散的程度．
\par
样本方差依赖于样本的选取，带有随机性．如果样本是随机抽取的，当样本容
\par
量较大时，样本方差是总体方差的估计．
\par
对于前面“思考”中的问题，由获得的样本数据，可以计算出甲、乙射击成绩
\par
的方差：
\par
一[（7一8）2十（8一8）2十，十（9一8）2〕一0．6，
\par
1
\par
1
\par
5乙
\par
一[（6一8）2十（8一8）2十，十（9一8）2〕一1.4．
\par
由于《<0，可以估计甲的射击成绩比乙更稳定，故可选择甲参加运动会．
\end{ocrdraft}
\end{bankitem}
